\subsection{Rectangular guards and unavoidable tree size}
\label{app:core-rectangles}

Repeated occurrences can be necessary in exponentially large numbers.
The obstruction is visible in the canonical source coefficients: one
occurrence separates its filters by owner, even when each filter has
many prime guards. We first record this rectangular structure.

Fix distinct owners $A,B,C$, put $N=I\setminus\{B\}$, and choose
disjoint passive variable blocks $\mathcal A,\mathcal C$ and a marked
source $j$ outside them. For $U\subseteq\mathcal A$, $W\subseteq\mathcal C$,
let $v_{U,W}$ give $\{B\}$ to $j$, $\{A\}$ to $U$, $\{C\}$ to $W$,
and $N$ to every other input. Define the source matrix
\[
 M_t(U,W)=[B\in t(v_{U,W})].
\]
A \emph{rectangle} is a Cartesian product of a family of $U$-sets
and a family of $W$-sets; it is positive if every entry in it is $1$.

\begin{lemma}[Rectangular source contributions]\label{core:lem:source-rectangles}
For any constant-free choice/union word $t$, the positive entries of
$M_t$ are a union of at most $\#_j(t)$ positive rectangles, where
$\#_j(t)$ counts occurrences of $j$. The rectangles can be chosen
downward closed in both coordinates. In particular, if $m$ positive
entries have the property that every two have at least one negative
crossed entry, then $\#_j(t)\geq m$.
\end{lemma}
\begin{proof}
Only $j$ supplies $B$. Fix one occurrence and follow its path.
At an $A$-gate it survives precisely when the sibling's $A$-profile
is false, a condition only on $U$. At a $C$-gate the condition only
involves $W$. A $B$-gate imposes no filter; any other owner's singleton
cannot occur on these inputs. Union imposes no filter either.
Thus this occurrence's contribution is a product of two antitone
predicates. Source decomposition gives the stated union. A positive
rectangle containing two selected entries would contain both crossed
entries, proving the last assertion.
\end{proof}

For distinct inputs $x,a_0,b_0,\ldots,a_{d-1},b_{d-1}$, define
\begin{equation}\label{core:eq:iterated-mixed-circuit}
 T_0=x,\qquad
 T_{i+1}=(a_i\star_A T_i)\star_B(T_i\star_C b_i).
\end{equation}
The two references to $T_i$ may share its output in a circuit.

\begin{theorem}[Exponential size throughout a core lower interval]
\label{core:thm:robust-tree-size}
For every $d\geq0$ and $K\geq3$, any constant-free choice/union word
$s$ on these inputs with $s(w)\subseteq T_d(w)$ at every input satisfies
\[
 \#_x(s)\geq2^d,\qquad
 \#_{a_i}(s),\#_{b_i}(s)\geq2^{d-i-1}\quad(0\leq i<d).
\]
Consequently it has at least $3\cdot2^d-2$ leaves. The unfolding of
\eqref{core:eq:iterated-mixed-circuit} attains all these bounds, so
$3\cdot2^d-2$ is the exact minimum even among all subordinate words.
The unfolded tree has $3(2^d-1)$ binary gates, whereas the shared circuit has $3d$
choice gates and $2d+1$ input nodes.
\end{theorem}
\begin{proof}
The profiles obtained directly from the recurrence are
\[
 \Phi_A^{T_d}=x\prod_{i<d}b_i,\qquad
 \Phi_C^{T_d}=x\prod_{i<d}a_i,\qquad
 \Phi_B^{T_d}=x\prod_{i<d}(a_i\vee b_i).
\]
We first show that every subordinate $s$ has exactly this $B$-profile.
All positive profile bits of $T_d$ persist in $s$ by nonemptiness.
Put $P_A=\{x,b_i:i<d\}$ and $P_C=\{x,a_i:i<d\}$.
Their intersection is $\{x\}$, so common-source compatibility forces
every $B$-winning coalition of $s$ to contain $x$.

Suppose such a coalition $E$ misses both members of pair $t$.
Set $w_x=\{B\}$ and
\[
 w_{a_i}=\begin{cases}\{C\},&a_i\in E,\\\{A\},&a_i\notin E,\end{cases}
 \qquad
 w_{b_i}=\begin{cases}\{A\},&b_i\in E,\\\{C\},&b_i\notin E.\end{cases}
\]
Stage $t$ gives $T_{t+1}(w)=\{A,C\}$, and all later stages preserve
this value. Hence $s(w)\subseteq\{A,C\}$.
Apply \cref{core:thm:one-bound-profiles} to the positive coalitions
$P_A,E,P_C$ of owners $A,B,C$, with full coalitions for additional
owners and output bound $\{A,C\}$. Their common intersection is
$\{x\}$, whose input is outside the bound. Moreover
\[
 E_A(w)\cap P_C\cap E=\varnothing,
 \qquad E_C(w)\cap P_A\cap E=\varnothing.
\]
Neither permitted absorber is possible, a contradiction. Thus the
$B$-winning coalitions are exactly those containing $x$ and at least
one member of every pair.

We use a separator observation. If
$\Phi_B^s(H)=1$, $\Phi_B^s(H\setminus\{j\})=0$, and a probe gives
$j$ the menu $\{B\}$ while every foreign atomic set misses $H$,
then its output contains $B$. Descend through children retaining the
positive/negative profile pair. At a $B$-gate choose a separating
child. At any foreign gate both children are $H$-positive; its
sibling cannot be its owner's singleton, by common-source compatibility
with $H$. A separating child therefore retains its $B$ contribution.
At a union root choose a separating child; union preserves its $B$.
The selected leaf must be $j$, proving the observation.

For source $x$, restrict $a_i$ to $\{A\},N$ and $b_i$ to $\{C\},N$.
Writing $U=\{i:w_{a_i}=\{A\}\}$ and
$W=\{i:w_{b_i}=\{C\}\}$, the recurrence gives
\begin{equation}\label{core:eq:disjointness-source}
 B\in T_d(v_{U,W})\quad\Longleftrightarrow\quad U\cap W=\varnothing.
\end{equation}
For each $S\subseteq\{0,\ldots,d-1\}$ take $(U,W)=(S,S^c)$.
The coalition consisting of $x$ and the neutral partner of each pair
is $B$-winning, and removing $x$ makes it losing. The separator
observation forces $B$ in $s$ at all these $2^d$ probes.
For distinct $S,R$, at least one cross has nonempty intersection,
so \eqref{core:eq:disjointness-source} and $s\leq T_d$ make that cross
negative for $s$. \Cref{core:lem:source-rectangles} gives
$\#_x(s)\geq2^d$.

For source $a_i$, set all inputs of $T_i$ and $b_i$ to $N$, and
$a_i=\{B\}$. Then $T_i=N$ and $T_{i+1}=I$. Vary only the higher
pairs as before. Their source matrix is the same disjointness matrix
on $d-i-1$ pairs. At each positive complementary pair, choose $H$
to contain $x$, one variable from each lower pair, $a_i$, and the
neutral partner of each higher pair. Exclude $b_i$ from $H$.
Then $H$ is $B$-winning, $H\setminus\{a_i\}$ is losing, and all
foreign atomic sets miss $H$. The same separator and rectangle
arguments give $\#_{a_i}(s)\geq2^{d-i-1}$. Interchanging $a_i,b_i$
gives the other bound.

Summing yields
$2^d+2\sum_{i<d}2^{d-i-1}=3\cdot2^d-2$ leaves. The displayed
unfolding has exactly these multiplicities. A binary tree has one
fewer gate than leaves, while each shared stage adds three gates.
\end{proof}

In particular, no bound depending only on $K$ limits the occurrences
of a chosen source in subordinate witnesses of positive probes.
Taking unions of words with that fixed source bound does not repair
the failure: every summand of a union equal to $T_d$ is subordinate.

\begin{corollary}[The reset bound survives the full core]
\label{core:cor:reset-ambient-tree-bound}
Let $R_0=c\del z$ and
$R_{i+1}=(R_i\del p_i)\del((R_i\del n_i)\del z)$.
For $K\geq3$, every core choice/union word inducing $R_d$ has at
least $2^d$ occurrences of $c$.
\end{corollary}
\begin{proof}
Set $c=\{B\}$ and $z=N$, and restrict $p_i$ to $\{A\},N$ and
$n_i$ to $\{C\},N$. No other input supplies $B$.
The selection formula in \cref{reset:prop:large-circuit-table} then gives
\[
 B\in R_d\quad\Longleftrightarrow\quad
 \forall i<d:\ p_i\text{ is ambiguous or }n_i\text{ is ambiguous}.
\]
This is exactly the disjointness matrix
\eqref{core:eq:disjointness-source}, whose $2^d$ complementary
pairs force the occurrence bound by \cref{core:lem:source-rectangles}.
\end{proof}
